\chapter{Plano de Projeto e Viabilidade}
\label{chap:plano_projeto}

\section{Termo de Abertura do Projeto (Project Charter)}
Este documento estabelece formalmente a autorização e o escopo estratégico para o desenvolvimento da plataforma de software, definindo objetivos corporativos, partes interessadas e alinhamento de mercado.

\subsection{Visão do Produto e Declaração de Posicionamento}
\textbf{Para} organizações de saúde e clínicas médicas, \textbf{cujo} desafio reside na descentralização operacional e fragmentação de prontuários, \textbf{o} Sistema de Gestão de Clínica (SGC) \textbf{é uma} plataforma corporativa de prontuário eletrônico e gestão operacional integrada \textbf{que} automatiza agendamentos, histórico clínico e faturamento em tempo real, \textbf{diferente de} sistemas legados engessados e planilhas desconexas, \textbf{nosso produto} oferece arquitetura modular assíncrona com soberania total de dados e conformidade estrita com a LGPD desde a concepção.

\subsection{Mapeamento de Partes Interessadas (Matriz Poder $\times$ Interesse)}
A gestão de stakeholders segue a Matriz de Mendelow para orientar as estratégias de comunicação e engajamento:
\begin{itemize}
    \item \textbf{Gerenciar de Perto (Alto Poder, Alto Interesse):} Diretores da Clínica, Corpo Clínico Chefe e Arquiteto de Software;
    \item \textbf{Manter Satisfeito (Alto Poder, Baixo Interesse):} Encarregado de Proteção de Dados (DPO) e Auditores Externos;
    \item \textbf{Manter Informado (Baixo Poder, Alto Interesse):} Recepcionistas, Atendentes e Pacientes;
    \item \textbf{Monitorar (Baixo Poder, Baixo Interesse):} Prestadores de serviços periféricos de manutenção física.
\end{itemize}

\section{Definição do Problema e Cenário Atual}
Apresentação detalhada das ineficiências, gargalos e riscos identificados no cenário atual que justificam a construção do software.

\section{Objetivos do Sistema}
\subsection{Objetivo Geral}
Desenvolver uma solução de software corporativa, segura e escalável para unificar a gestão assistencial e administrativa.

\subsection{Objetivos Específicos}
\begin{itemize}
    \item Automatizar os processos operacionais críticos da recepção e consultórios;
    \item Garantir rastreabilidade integral dos registros e auditoria imutável de ações clínicas;
    \item Assegurar conformidade técnica com requisitos de segurança e legislação aplicável (LGPD).
\end{itemize}

\section{Estudo de Viabilidade e Soluções Alternativas}
\begin{table}[htbp]
\caption{Matriz Comparativa de Abordagens de Solução}
\label{tab:comparativo_solucoes}
\centering
\small
\begin{tabularx}{\textwidth}{|l|Y|Y|l|}
\hline
\rowcolor{tableheaderbeige}
\textbf{Abordagem} & \textbf{Vantagens} & \textbf{Desvantagens} & \textbf{Veredito} \\ \hline
Desenvolvimento Próprio & Aderência total ao negócio, controle do código e soberania LGPD. & Maior tempo inicial de construção. & \textbf{Adotada} \\ \hline
Software as a Service (SaaS) & Implantação imediata e suporte do fornecedor. & Custo recorrente elevado, dependência de terceiros. & Descartada \\ \hline
Customização Open-source & Custo zero de licença e base pré-existente. & Esforço de adaptação equivalente a refazer o núcleo. & Descartada \\ \hline
\end{tabularx}
\end{table}

\subsection{Solução Adotada e Justificativa Técnica (Rationale)}
Optou-se pelo desenvolvimento próprio do sistema para assegurar soberania sobre os dados sensíveis dos pacientes, eliminar custos de licenciamento recorrentes por assento e permitir total controle sobre a arquitetura de segurança, garantindo conformidade com a LGPD desde a concepção.

\section{Especificação de Componentes e Infraestrutura}
\subsection{Hardware e Periféricos}
Especificação dos nós de produção, estações de trabalho e dispositivos móveis necessários para operação.

\subsection{Stack de Software}
Linguagens, frameworks, bibliotecas e sistemas gerenciadores de banco de dados adotados.

\subsection{Redes e Hospedagem}
Topologia de rede local, links redundantes de comunicação externa e provedor de computação em nuvem.

\section{Organização da Equipe de Desenvolvimento}
\begin{table}[htbp]
\caption{Alocação de Responsabilidades da Equipe (Matriz RACI Sintética)}
\label{tab:equipe_raci}
\centering
\small
\begin{tabularx}{\textwidth}{|l|l|Y|}
\hline
\rowcolor{tableheaderbeige}
\textbf{Perfil / Função} & \textbf{Responsável} & \textbf{Responsabilidades Principais} \\ \hline
Arquiteto de Software & Engenheiro de Software Sênior & Decisões estruturais, modelos 4+1 e ADRs. \\ \hline
Engenheiro de Requisitos & Analista de Negócios / Requisitos & Elicitação, especificação formal e validação. \\ \hline
Desenvolvedores Front/Back & Equipe de Desenvolvimento & Construção e testes unitários/integrados. \\ \hline
Engenheiro de DevOps/QA & Especialista em Infra/Testes & CI/CD, ambientes, pipelines de teste e deploy. \\ \hline
\end{tabularx}
\end{table}

\section{Gerenciamento de Tempo e Cronograma}
\subsection{Estrutura Analítica do Projeto (WBS / EAP)}
Decomposição hierárquica do escopo em pacotes de trabalho mensuráveis. A \autoref{fig:proj_wbs} ilustra a decomposição em árvore das entregas.

% PLACEHOLDER STARUML v7.0: Estrutura Analítica do Projeto (WBS / EAP)
% Ferramenta: StarUML v7.0 -> Model -> Add Diagram -> Class Diagram (Hierarquia em Árvore)
% Arquivo: Imagens/proj_wbs_escopo.png (Exportar em PNG 300 DPI ou PDF vetorial)
% Descomente a linha abaixo apos salvar o arquivo no diretorio Imagens/:
% \incluirdiagrama{Imagens/proj_wbs_escopo.png}{Estrutura Analítica do Projeto (WBS / EAP)}{fig:proj_wbs}
\begin{figure}[htbp]
    \centering
    \fbox{\parbox{0.9\textwidth}{\centering\vspace{1cm}%
        \textbf{[Estrutura Analítica do Projeto (WBS) -- StarUML v7.0]}\\%
        \textit{Modelar no StarUML v7.0 e exportar para: \texttt{Imagens/proj\_wbs\_escopo.png}}\\%
        \small Menu: \texttt{File -> Export Diagram as -> PNG...} (Fundo branco, alta resolução 300 DPI)\\%
        \vspace{0.3cm}%
        \footnotesize Para ativar a imagem real, descomente no arquivo \texttt{.tex}:\\
        \texttt{\textbackslash incluirdiagrama\{Imagens/proj\_wbs\_escopo.png\}\{Estrutura Analítica do Projeto (WBS)\}\{fig:proj_wbs\}}%
    \vspace{1cm}}}
    \caption{Estrutura Analítica do Projeto (Placeholder StarUML v7.0)}
    \label{fig:proj_wbs}
\end{figure}

\subsection{Cronograma Sintético e Marcos de Entrega}
\begin{table}[htbp]
\caption{Cronograma Macro e Marcos Formais de Entrega}
\label{tab:marcos_entrega}
\centering
\small
\begin{tabularx}{\textwidth}{|l|l|Y|l|}
\hline
\rowcolor{tableheaderbeige}
\textbf{Marco} & \textbf{Período} & \textbf{Entregáveis Consolidados} & \textbf{Status} \\ \hline
M1: Planejamento & Semanas 1 a 2 & Plano de Projeto, WBS e Estimativas. & \badgeconcluido \\ \hline
M2: Requisitos & Semanas 2 a 5 & Especificação de Requisitos, Casos de Uso e Protótipos. & \badgeconcluido \\ \hline
M3: Arquitetura/Design & Semanas 5 a 7 & Documento de Arquitetura, Modelo de Dados e Diagramas UML. & \badgeconcluido \\ \hline
M4: Construção & Semanas 7 a 12 & Módulos implementados e integrados. & \badgeemprogresso \\ \hline
M5: Homologação & Semanas 13 a 14 & Bateria de testes, deploy e validação final. & \badgenaoiciada \\ \hline
\end{tabularx}
\end{table}

\subsection{Diagrama de Gantt do Projeto}
A \autoref{fig:proj_gantt} apresenta a distribuição temporal das frentes de trabalho, dependências de precedência e o caminho crítico (CPM) implementados nativamente via \texttt{pgfgantt}.

\begin{figure}[htbp]
    \centering
    \resizebox{\textwidth}{!}{
    \begin{ganttchart}[
        % Configuração de Grade
        vgrid={*1{draw=gray!25, dashed}},
        hgrid={*1{draw=gray!20}},
        x unit=0.8cm,
        y unit chart=0.60cm,
        y unit title=0.55cm,
        title/.append style={fill=white, draw=black, thick},
        title label font=\bfseries\footnotesize,
        % Barras Padrão (Frentes Paralelas / Secundárias)
        bar/.append style={fill=blue!20!purple!25, draw=blue!50!purple!70, rounded corners=1pt},
        bar height=0.50,
        bar label font=\scriptsize,
        % Grupos de Fases
        group/.append style={fill=gray!25, draw=black!80, rounded corners=1.5pt},
        group height=0.25,
        group label font=\bfseries\scriptsize,
        % Marcos Formais de Entrega
        milestone/.append style={shape=diamond, fill=black, draw=black, scale=0.9},
        milestone label font=\bfseries\scriptsize
    ]{1}{12} % Definição do intervalo total de semanas (ex: 1 a 12)
        
        % 1. Cabeçalho de Meses (Número de semanas por mês)
        \gantttitle{Mês 1}{3}
        \gantttitle{Mês 2}{5}
        \gantttitle{Mês 3}{4} \\
        
        % 2. Cabeçalho de Semanas
        \gantttitle{S1}{1} \gantttitle{S2}{1} \gantttitle{S3}{1}
        \gantttitle{S4}{1} \gantttitle{S5}{1} \gantttitle{S6}{1} \gantttitle{S7}{1} \gantttitle{S8}{1}
        \gantttitle{S9}{1} \gantttitle{S10}{1} \gantttitle{S11}{1} \gantttitle{S12}{1} \\
        
        % FASE 1
        \ganttgroup{Fase 1: Concepção \& Infraestrutura}{1}{2} \\
        % Para tarefa no Caminho Crítico, sobrescrever estilo da barra:
        \ganttbar[bar/.append style={fill=magenta!35!purple!55, draw=magenta!70!black}]{1.1 Atividade Crítica A}{1}{2} \\
        % Para tarefa paralela, usar o estilo padrão:
        \ganttbar{1.2 Atividade Paralela B}{1}{2} \\
        
        % FASE 2
        \ganttgroup{Fase 2: Serviços Núcleo (MVP)}{2}{6} \\
        \ganttbar[bar/.append style={fill=magenta!35!purple!55, draw=magenta!70!black}]{2.1 Atividade Crítica C}{2}{4} \\
        \ganttbar{2.2 Atividade Paralela D}{3}{4} \\
        \ganttbar[bar/.append style={fill=magenta!35!purple!55, draw=magenta!70!black}]{2.3 Atividade Crítica E}{5}{6} \\
        \ganttmilestone{Parte 1: Entrega Parcial (Data)}{6} \\
        
        % FASE 3
        \ganttgroup{Fase 3: Construção \& Integração}{7}{9} \\
        \ganttbar[bar/.append style={fill=magenta!35!purple!55, draw=magenta!70!black}]{3.1 Atividade Crítica F}{7}{8} \\
        \ganttbar{3.2 Atividade Paralela G}{8}{9} \\
        
        % FASE 4
        \ganttgroup{Fase 4: Homologação \& Entrega}{10}{12} \\
        \ganttbar[bar/.append style={fill=magenta!35!purple!55, draw=magenta!70!black}]{4.1 Testes de Carga e Validação}{10}{11} \\
        \ganttbar{4.2 Documentação Final}{11}{12} \\
        \ganttmilestone{Parte 2: Entrega Final (Data)}{12}
    \end{ganttchart}
    }
    \caption{Diagrama de Gantt com mapeamento do Caminho Crítico (em destaque), marcos de entrega e frentes paralelas de desenvolvimento.}
    \label{fig:proj_gantt}
\end{figure}

\section{Análise e Gerenciamento de Riscos}
\begin{table}[htbp]
\caption{Registro de Riscos do Projeto}
\label{tab:registro_riscos}
\centering
\small
\begin{tabularx}{\textwidth}{|l|Y|c|c|Y|}
\hline
\rowcolor{tableheaderbeige}
\textbf{ID} & \textbf{Descrição do Risco} & \textbf{Prob.} & \textbf{Imp.} & \textbf{Ação de Mitigação / Prevenção} \\ \hline
R01 & Desalinhamento de requisitos de negócio & Média & Alto & Prototipação rápida e validação formal. \\ \hline
R02 & Falhas de concorrência em horários e estoque & Média & Alto & Bloqueio pessimista / isolamento transacional ACID. \\ \hline
R03 & Vazamento de dados pessoais/sensíveis & Baixa & Crítico & Criptografia, RBAC granular e logs de auditoria. \\ \hline
\end{tabularx}
\end{table}

\section{Estimativas de Custos e Orçamento}
Projeção orçamentária considerando infraestrutura em nuvem, ferramentas de apoio, licenças e mão de obra alocada.
